\section{Baseline NL}
The Baseline NL module provides a reference framework for evaluating Large Language Models (LLMs) through direct interaction via natural language prompts.
This approach assesses the model’s intrinsic reasoning and knowledge by prompting it with natural language (NL) questions, optionally enriched with structural and contextual information derived from a database schema or a specific query results.
All outputs are standardized in a unified FULL JSON format to facilitate downstream evaluation.

\subsection{General Procedure}
\begin{enumerate}
    \item  \textbf{Prompt Loading:} Natural language prompts are loaded from dataset-specific resources (\textbf{queries\_<dataset-name>.txt}). Each prompt corresponds to a query from the same dataset.
    \item \textbf{LLM interaction:} Prompts are processed through a \textbf{LangChain-based pipeline}, which integrates:
    \begin{itemize}
        \item The general task definition and LLM role specification.
        \item The corresponding database schema, extracted via DuckDB.
        \item Optional contextual data (e.g. query results for Model Context datasets).
        \item The natural language question itself
    \end{itemize}
    \item \textbf{Response Parsing and Storage:} The model's output is parsed using a \textbf{PyddanticOutputParser}, ensuring compliance with the required FULL JSON format. Each response is stored in an individual JSON file, forming a consistent source for evaluation.
\end{enumerate}


\subsection{Architecture:}
The architecture is composed of three main components:
\begin{enumerate}
    \item \textbf{Configuration Layer:} 
    \begin{itemize}
        \item API keys for the supported providers (\textbf{IBMWatsonx} and \textbf{Google Gemini}) are loaded form a \textbf{.env} file.
        \item The \textbf{Config\_Loader} selects the target LLM provider, and the appropriate wrapper is instantiated through \textbf{get\_llm\_wrapper()}.
        \item Each wrapper extends the \textbf{LLMBaseWrapper} class, implementing provider-specific logic for model initialization, token counting, and response handling.
    \end{itemize}
    \item \textbf{Chain Construction:}
    
    The \textbf{build\_lcel\_chain()} method constructs a \textbf{LangChain Expression Language (LCEL)} pipeline comprising:
    \begin{itemize}
        \item \textbf{Database Context Retrieval:} \textbf{get\_db\_context()} extracts schema information and, for Model Context datasets, executes the original SQL query using DuckDB database instance of corresponding dataset.
        \item \textbf{Prompt Composition:} \textbf{ChatPromptTemplate} merges two layers, a \textbf{SYSTEM\_PROMPT} defining the LLM's role and behavior, and a \textbf{HUMAN\_PROMPT} containing the NL question and formatting directives.
        \item \textbf{Structured Output Parsing:} The \textbf{PydanticOutputParser} enforces adherence to the predefined FULL JSON schema.
    \end{itemize}
    \item \textbf{Execution Flow}
    \begin{itemize}
        \item For each dataset, SQL queries and corresponding NL prompts are loaded.
        \item The \textbf{(prompt, query)} pairs  are processed through \textbf{chain.invoke()}.
        \item Responses are parsed, validated, and augmented with the metadata such as execution time and token usage.
        \item The resulting files are serialized via \textbf{save\_baseline\_to\_json()} for evaluation purposes.
    \end{itemize}
\end{enumerate}

